%%==============================================================================
%% sections/01_introduction.tex
%%==============================================================================
\section{Introduction}
\label{sec:intro}

The remarkable progress of large language models
\citep{brown2020gpt3,touvron2023llama,openai2023gpt4} has been driven primarily
by scale. However, this scaling paradigm is increasingly out of reach for the
broader research community: training a 70B-parameter model from scratch
requires tens of millions of dollars of compute.

In this work we ask: \emph{can careful data engineering and algorithmic design
close the gap between compact and frontier LLMs on reasoning tasks?} We answer
this question affirmatively with BJTU-LLM, and identify three ingredients that
prove crucial:

\begin{contributions}
\begin{enumerate}
    \item \textbf{Data.} A curriculum-based synthetic-data pipeline
          (\S\ref{sec:method-data}) that generates 12B tokens of high-quality
          reasoning traces spanning mathematics, code, and logic.
    \item \textbf{Architecture.} A hardware-aware sparse-attention operator
          (\S\ref{sec:method-arch}) that reduces long-context latency by
          $2.3\times$ on NVIDIA H800 GPUs without accuracy loss.
    \item \textbf{Alignment.} A two-stage RLHF procedure with process-level
          rewards (\S\ref{sec:method-rlhf}) that yields consistent gains on
          multi-step reasoning benchmarks.
\end{enumerate}
\end{contributions}

Extensive experiments (\S\ref{sec:experiments}) demonstrate that BJTU-LLM-7B
matches or exceeds Llama-3-70B on four out of six reasoning benchmarks, while
requiring only $3.1$\% of the training FLOPs.
